\documentclass{projectreport}

\title{Project title}
\date{March 20, 2022}

\coursecode{COMP307}
\y{III} % Year
\sem{II} % Semester/Part
\group{Computer Engineering} % Computer Science/Engineering

\member{Name 1}{Roll}
\member{Name 2}{Roll}
\member{Name 3}{Roll}
%\member{Name 4}{Roll}
%\member{Name 5}{Roll}

\supervisor{Supervisor's name}{Supervisotr's academic title}{Supervisor's department}


\coordinator{Coordinator Name}

% For abbreviations/Acronyms
\usepackage{glossaries}
\loadglsentries{glossaries}
\makenoidxglossaries

\begin{document}
	\maketitle
	
	\makebonafidecertificate
	
	\begin{acknowledgements}
		....		
	\end{acknowledgements}
	
	
	
	
	\begin{abstract}
		It should include summary of your project work. It should answer three main questions. What did you do? Why did you do? And how did you do it? The abstract should cover exact problem and how your work is going to address those issues. Abstract provides description about the project in minimum possible words. Abstract should contain 150-200 words as follows:
		
		\begin{itemize}
			\item Background (1 to 2 sentences)
			\item Purpose/Aim (1 to 2 sentences)
			\item Procedure and Method/Tools (1 to 2 sentences)
			\item Result (1 to 2 sentences)
			\item Conclusion (1 sentence)
			\item Recommendation (1 sentence)
		\end{itemize}
		
	\end{abstract}

	\begin{keywords}
		Keyword1, Keyword2
	\end{keywords}

	\pagenumbering{roman}
	\setcounter{page}{1}
		
	\tableofcontents
	\listoffigures
	\listoftables
	
	\clearpage
	\printnoidxglossary[ title=Abbreviations]
	
	
	\chapter{Introduction}
	\pagenumbering{arabic}
	\setcounter{page}{1}
	This chapter introduces the developed system. It is an elaborated form of Abstract, and is usually divided into subsections, such as background, objectives, motivation, and significance of the project.
	
	\section{Background}
	It should include basic information in the field of project work. Recent trends and development on the area of the project should be briefly discussed in this topic. It should contain at least 2-3 paragraphs, having 200-400 words.
	It should include:
	
	\begin{itemize}
		\item Recent development in the field
		\item Drawbacks and significance of existing works
	\end{itemize}
	
	\section{Objectives}
	What is the basic purpose of your project? List the objectives of your project work in bullets (not exceeding 4).
	
	\section{Motivation and Significance}
	Motivation to choose the topic, importance and contribution of your project work is included in this section.
	
	\begin{itemize}
		\item Why did you choose this particular topic as your project?
		\item How does the work address drawback of the existing systems?
		\item How is it different from the existing works?
	\end{itemize}

	It consists of a short paragraph about features of your project work. A brief introduction of features that project is going to address, helps to make the project proposal more robust.
	
	\chapter{Related Works}
	
	It focuses on the discussion of similar types of other tasks/projects that have been performed earlier. It should include recent projects, works, websites, etc. as reference. Referencing should be in APA format. Below are some examples:
	
	According to \cite{doe2012}, Management Systems should include components of ...
	
	
	Thapa and Shrestha (2010) have done similar projects where they have implemented ...
	
	The Minesweeper \cite{rai2012} has features like ...
	
	Similar kinds of project work were performed earlier which have properties ...
	
	
	
	In writing this section, your purpose is to convey to your readers what knowledge and ideas have been established on a topic, and what their strengths and weaknesses are \cite{taylor2008literature}. It should also include previous works, comparison between works, drawbacks and current works.
	
	
	\chapter{ Design and Implementation}
	It includes the explanation of the sequential procedure that you performed during the project work. It may include algorithms, flowcharts, System Diagrams (Use-case, \gls{er}, Class, Activity etc.) and other analysis and design which give general idea about how you approach different development procedures during the project work. It may contain a few paragraphs and diagrams which may describe your work. 
	
	
	\section{System Requirement Specifications}
	
	This section specifies the requirements of the developed system. It may include software specifications and hardware specifications.
	
	\subsection{Software Specifications}
	This section presents the description of the developed software system. It should include the functional and non-functional requirements of the software. It may also include software dependencies.
	
	Figure \ref{fig:usecase}\footnote{https://en.wikipedia.org/wiki/Use\_case} is an example of a use-case diagram. Table \ref{table:sample} is an example of a table.
	
	\begin{figure}
		\centering
		\includegraphics[width=\textwidth]{usecase}
		\caption{A sample use-case diagram}
		\label{fig:usecase}
	\end{figure}

	\begin{table}[h!]
		\caption{A sample table}
		\label{table:sample}
		\centering
		\begin{tabular}{|l|l|l|}
			\hline
			\textbf{Header} & \textbf{Header} & \textbf{Header}\\ \hline
			Content & Content & Content \\ \hline
		\end{tabular}
	\end{table}
	
	\subsection{Hardware Specifications}
	It contains description of any additional hardware that may be required for the project. It also gives information about required minimum configuration of the system to run the system smoothly.
	
	
	\chapter{Discussion on the achievements}
	This chapter should focus on challenges that you have experienced during your work, and discuss any probable deviations from the objectives, negative findings that were implemented.
	
	Any measurable results and comparisons that were obtained from the work could be included in this section.
	
	\section{Features}
	Include/list out information about the new/vital features that has been implemented in project development, with brief discussion.
	
	\chapter{Conclusion and Recommendation}
	Summarize your achieved/unachieved goals with valid reasoning. The conclusion should not contradict with the objectives of the project.
	
	\section{Limitations}
	Include limitation of your work.
	
	
	\section{Future Enhancement}
	Include possible enhancement of this work
	
	\bibliography{report}
	
	
\end{document}